\documentclass{article}
\usepackage[left=2cm, right=2cm, top=2cm, bottom=2cm]{geometry}
\usepackage{graphicx}
\usepackage{amsmath}
\usepackage{amssymb}
\usepackage{amsthm}
\usepackage{fancyhdr}
\usepackage{verbatim}
\usepackage{listings}
\usepackage{xcolor}
\usepackage{pdfpages}
\usepackage{cancel}
\usepackage{physics}
\usepackage{esint}

\lstset{
    language=C++,
    basicstyle=\ttfamily\footnotesize,
    keywordstyle=\color{blue},
    commentstyle=\color{green},
    stringstyle=\color{red},
    numbers=left,
    numberstyle=\tiny\color{gray},
    frame=single,
    breaklines=true,
    captionpos=b,
}


\title{MTH 553: Lecture \# 31}
\author{Cliff Sun}

\newtheorem{theorem}{Theorem}[section]
\newtheorem{lemma}[theorem]{Lemma}
\newtheorem{definition}[theorem]{Definition}
\newtheorem{conjecture}[theorem]{Conjecture}
\newtheorem{proposition}[theorem]{Proposition}
\newtheorem{corollary}[theorem]{Corollary}
\newtheorem{one minute paper}[theorem]{One Minute Paper}

\pagestyle{fancy}
\lhead{\textbf{\thepage}\ \ \nouppercase{\rightmark}}
\chead{MTH 553: Lecture \# 31}
\rhead{Cliff Sun}

\begin{document}

\maketitle

\section*{Lecture Span}
\begin{enumerate}
    \item Fundamental solutions
\end{enumerate}

If $f \in C_0^2(\mathbb{R}^n)$, then $u = k \star f$ solves $-\triangle u = f$ classically. 

\begin{proof}
    \begin{equation}
        u(x) = \int_{\mathbb{R}^n}K(y)f(x-y)dy
    \end{equation}
    Fix $x$, then 
    \begin{align*}
        u_{x_i}(x) &= \lim_{h \rightarrow 0}\frac{u(x + he_i) - u(x)}{h} \\
        &= \lim_{h\rightarrow 0}\int K(y)\frac{f(x-y + he_i) - f(x-y)}{h}dy \\
        &= \int K(y) f_{x_i}(x-y)dy
    \end{align*}
    Similarly,
    \begin{equation}
        u_{x_i,x_j} = \int K(y)f_{x_ix_j}(x-y)dy
    \end{equation}
    Therefore, $u_{x_i,x_j}$ is continuous and $u \in C^2(\mathbb{R}^n)$. 
\end{proof}

Since $u$ is $C^2$ smooth and $-\triangle u = f$ weakly, then conclude $-\triangle u = f$ classically. 

\section*{Potential Theory}

We have solved Poisson on $\mathbb{R}^n$. Then, solve on bounded domains using Greens functions but rather than just eigenfunctions. 

\begin{theorem}
    \textbf{(Representation Formula for Poisson's Equation on Bounded Domain)} If $\Omega \subset \mathbb{R}^n$ and $n \geq 2$ is smooth and bounded, $u \in C^2(\overline{\Omega})$ and $x \in \Omega$. 
    Then, 
    \begin{equation}
            u(x) = \int_{\Omega}K(x - y)(-\triangle u)(y)dy + \int_{\partial \Omega}\left( -u(y)\partial_{v_y}K(x-y) + K(x-y)\partial_{v_y}u \right)dS    
    \end{equation}
\end{theorem}

\begin{proof}
    By arguing like last time (when we showed that $-\triangle K = \delta_0$), we obtain that 
    \begin{equation}
        \int_{B(x,\epsilon)}K(x-y)\triangle u(y)dy \rightarrow 0
    \end{equation}
    As $\epsilon\rightarrow 0$. Moreover, 
    \begin{equation}
        \int_{\Omega / B(x,\epsilon)}K(x-y)\triangle u(y)dy \rightarrow -u(x) + \int_{\partial \Omega}K(x-y)\partial_{v_y}u(y)dS(y) - \int_{\partial \Omega}\partial_{v_y}K(x-y)u(y)dS(y)
    \end{equation}
    As $\epsilon\rightarrow 0$. Previously, we can ignore the terms on $\partial \Omega$ since $\Omega$ had compact support. But now, we have to keep them... Adding equations (4) and (5) and moving some terms around proves the theorem.
\end{proof}

\textbf{What do the terms mean?}: THe first term is the domain potential. The second term is called the "double layer potential". The last term is called the "single layer potential". 

\subsection*{Electrostatics in $\mathbb{R}^3$}

Consider a point charge in $\mathbb{R}^3$ at $z$. Then, the electric field at $x$ has magnitude of $1/(x-z)^2$. Therefore, 
\begin{equation}
    E(x) = \frac{1}{|x-z|^2}\frac{x-z}{|x-z|}
\end{equation}

Note that 

\begin{equation}
    E = -\nabla u, \quad u(x) = \frac{1}{|x-z|}
\end{equation}

The domain potential is 

\begin{equation}
    u(x) = 4\pi \int_{\mathbb{R}^3}K(x-y)\rho(y)dy
\end{equation}

Where $\rho(y)$ is the charge distribution. We think of the right hand side of the Poisson's equation as a charge distribution. For the single layer potential, 

\begin{equation}
    4\pi \int_{\partial \Omega}K(x-y)\sigma(y)dy
\end{equation}

Here, $\sigma(y)$ is the surface charge density $= \partial_{v_y}u$ on $\partial \Omega$. Lastly, the double layer potential 

\begin{equation}
    4\pi \int_{\partial \Omega}\frac{\partial K}{\partial v_y} \mu(y)dS(y)
\end{equation}

Where $\mu$ is the dipole moment of the charge. Note, $\partial_{v_y}K$ is the derivative of $K$ in the direction of the normal vector. This is a double layer potential the derivative of $K$ is non zero if there is a change in the charge in the boundary. 

We also can't create a solution formula from the representation formula for the Poisson Equation in $\Omega$. This is because we can only specify the normal derivative or $u$, but not both. 

\end{document}